\subsection{Evaluation Setting}

We evaluate the benchmark as a multilingual document retrieval task. After dataset construction, the released \texttt{corpus}, \texttt{queries}, and \texttt{qrels} splits are loaded directly from the Hugging Face dataset repository and wrapped as a custom MTEB retrieval task. Query languages are inferred from the released metadata and mapped into the language labels expected by the evaluation framework so that the evaluation remains multilingual without requiring a manually maintained task manifest. The default reported setup is the \texttt{multilingual} variant; a stricter \texttt{cross\_language} variant is also available and keeps only other-language positives for each query.

\subsection{Models and Metrics}

The current evaluation pipeline compares multilingual sentence-embedding models under a common retrieval protocol. The default baseline set includes \texttt{sentence-transformers/paraphrase-multilingual-MiniLM-L12-v2}, \texttt{sentence-transformers/paraphrase-multilingual-mpnet-base-v2}, \texttt{intfloat/multilingual-e5-large}, and \texttt{BAAI/bge-m3}. We report standard ranking metrics, with Recall@10 treated as the main score. This is deliberate: each query can have several relevant document realizations in different retained languages, so the primary evaluation question is how much of that positive set the model retrieves in the top 10 rather than only whether one relevant document appears early. The saved benchmark artifacts additionally include Recall, MAP, nDCG, and same-language irrelevant-result share at \(k \in \{10,20,50,100\}\), plus language-specific \texttt{same\_language\_irrelevant\_share\_at\_100\_lang\_\{lang\}} diagnostics for German, English, Spanish, French, and Chinese queries. The compact paper table shows the 10 and 100 cutoffs rather than every stored diagnostic column.

\subsection{Implementation Notes}

Model artifacts are cached locally so repeated runs remain reproducible and independent of user-profile caches. The evaluation code writes structured result summaries to disk and generates comparison tables from saved result files rather than rerunning the full benchmark each time the paper tables need to be refreshed. The current pipeline relies on sentence-transformer model loading and the default backend device selection of the runtime environment.

\subsection{Results Placeholder}

This subsection is intentionally left without numeric results for now. Final model scores, ranking tables, and any language-wise breakdowns should be inserted here once the benchmark runs are finalized.

\begin{table}[t]
\centering
\caption{Placeholder for the main multilingual retrieval results.}
\label{tab:main-results-placeholder}
\begin{tabular}{lcccccccc}
\toprule
Model & Recall@10 & Recall@100 & MAP@10 & MAP@100 & MAP & nDCG@10 & nDCG@100 & SameLangIrr@100 \\
\midrule
% Results to be added.
\bottomrule
\end{tabular}
\end{table}

\subsection{Analysis Placeholder}

This subsection should later summarize the main empirical findings, including which encoder families transfer best to multilingual legal retrieval, whether smaller multilingual baselines remain competitive, and where the benchmark appears most challenging. Any language-specific or document-type-specific error analysis can be added here after the core results stabilize.
